\begin{lemma}
\label{editorial-source-lemma-mini-separability}
Let $k$ be a field of characteristic $p > 1$. Let $K/k$
be a field extension generated by $x_1, \ldots, x_{n + 1} \in K$ such that
\begin{enumerate}
\item $\{x_1, \ldots, x_n\}$ is a transcendence base of $K/k$,
\item for every $k$-linearly independent subset
$\{a_1, \ldots, a_m\}$ of $K$ the set
$\{a^p_1, \ldots, a_m^p\}$ is $k$-linearly independent.
\end{enumerate}
Then there is $1 \leq j \leq n+1$ such that
$\{ x_1, \ldots, \widehat{x}_j, \ldots, x_{n+1}\}$
is a separating transcendence base for $K / k$.
\end{lemma}

\begin{proof}
Retain the original generators $x_1,\ldots,x_{n+1}$.
Since $x_{n+1}$ is algebraic over $k(x_1,\ldots,x_n)$,
clearing the nonzero polynomial denominators in one of its
algebraic equations gives a nonzero
$F\in k[X_1,\ldots,X_{n+1}]$ with $F(x_1,\ldots,x_{n+1})=0$.
Choose such an $F$ of least total degree $D$. A nonzero
constant cannot vanish, so $D\geq1$. If $F=GH$ with both
factors of positive degree, evaluation in the field $K$
makes at least one factor zero. That factor has degree
less than $D$, a contradiction. Thus $F$ is irreducible.

Write the full polynomial as
$$
F=\sum_{\alpha\in A}\lambda_\alpha
X_1^{\alpha_1}\cdots X_{n+1}^{\alpha_{n+1}},
\qquad 0\neq\lambda_\alpha\in k,
$$
where $A$ is its finite set of distinct exponent tuples,
all of whose coordinates are integers at least zero.
Suppose every coordinate of every $\alpha\in A$ were
divisible by $p$. Put
$b_\alpha=x_1^{\alpha_1/p}\cdots x_{n+1}^{\alpha_{n+1}/p}$.
The family $(b_\alpha^p)_{\alpha\in A}$ is dependent,
since $\sum_{\alpha\in A}\lambda_\alpha b_\alpha^p=0$.
Hypothesis (2) also preserves independence of indexed
finite families: an independent family has distinct
entries, and Frobenius is injective on $K$, so its
$p$th powers still have distinct entries and hypothesis
(2) applies to the corresponding sets. Its contrapositive
therefore makes $(b_\alpha)_{\alpha\in A}$ dependent,
including the case where two evaluated monomials agree.
Choose $c_\alpha\in k$, not all zero, such that
$\sum_\alpha c_\alpha b_\alpha=0$. The polynomial
$$
G=\sum_{\alpha\in A}c_\alpha
X_1^{\alpha_1/p}\cdots X_{n+1}^{\alpha_{n+1}/p}
$$
is nonzero because its formal monomials are distinct.
It vanishes at the same original generators and has
degree at most $D/p<D$, a contradiction. Consequently
some exponent of some $X_i$ in $F$ is not divisible
by $p$, equivalently $\partial F/\partial X_i\neq0$.

Fix any such $i$, and retain the original coefficient
polynomials in
$$
F=\sum_{r=0}^h A_r(X_1,\ldots,\widehat X_i,\ldots,X_{n+1})X_i^r.
$$
There is $r\geq1$ with $p\nmid r$ and $A_r\neq0$.
Every nonzero $A_r$ with $r\geq1$ has total degree
at most $D-r<D$. Hence its value at the other original
generators is nonzero: otherwise that same coefficient
polynomial, viewed in all $n+1$ variables, would be a
smaller nonzero relation. In particular the polynomial
$$
P(T)=F(x_1,\ldots,x_{i-1},T,x_{i+1},\ldots,x_{n+1})
\in L[T],\qquad
L=k(x_1,\ldots,\widehat x_i,\ldots,x_{n+1}),
$$
is nonconstant and has $P(x_i)=0$. Thus $K=L(x_i)$
is finite algebraic over $L$. A transcendence basis of
$L/k$, chosen among its $n$ displayed generators by
Fields, Lemma \ref{fields-lemma-transcendence-degree},
is also a transcendence basis of $K/k$: algebraicity
in a tower is transitive. Its size must therefore be
$n$. All the displayed generators of $L$ are consequently
algebraically independent over $k$.

Evaluation thus identifies the polynomial ring
$B=k[X_1,\ldots,\widehat X_i,\ldots,X_{n+1}]$
with $k[x_1,\ldots,\widehat x_i,\ldots,x_{n+1}]$,
and identifies its fraction field with the original $L$.
The irreducible polynomial $F\in B[X_i]$ is primitive:
a nonunit common divisor of its coefficients would
factor $F$, with the other factor still of positive
$X_i$-degree. Gauss' lemma now makes the unchanged
polynomial $P(T)$ irreducible in $L[T]$.
Its derivative is the full sum
$$
P'(T)=\sum_{r=1}^h r A_r(x_1,\ldots,\widehat x_i,\ldots,x_{n+1})T^{r-1}.
$$
The previously chosen term is nonzero, so $P'\neq0$.
Irreducibility and $\deg P'<\deg P$ give
$\gcd(P,P')=1$. Hence $x_i$ is separable over $L$,
and $K/L$ is separable. This proves the assertion,
also when $n=0$ and the displayed basis of $L=k$ is empty.

The proof gives an exact criterion for every omitted
coordinate, not merely existence of one. For this same
minimal-degree $F$, omission of $x_j$ gives a separating
transcendence basis if and only if
$\partial F/\partial X_j\neq0$. Sufficiency was proved
for every such index. For necessity, suppose the other
$n$ elements form a separating basis and put
$L_j=k(x_1,\ldots,\widehat x_j,\ldots,x_{n+1})$.
The polynomial $F$ must depend on $X_j$, since otherwise
it would be a relation among that basis. Its coefficient
evaluation is injective, and the same primitive Gauss
argument makes its image $P_j(T)$ irreducible over $L_j$.
If $c_j$ is its actual leading coefficient and $M_j(T)$
is the monic minimal polynomial of $x_j$ over $L_j$,
then $P_j(T)=c_jM_j(T)$ with $c_j\neq0$. Separability
gives $M_j'\neq0$, so $P_j'=c_jM_j'\neq0$.
Injectivity of coefficient evaluation yields
$\partial F/\partial X_j\neq0$, as required.

There is also an exact relation-ring comparison for
these original generators. Evaluation induces
$$
k[X_1,\ldots,X_{n+1}]/(F)
\longrightarrow k[x_1,\ldots,x_{n+1}].
$$
It is surjective. If an original polynomial $H$ is in
its kernel, view $H$ in $B[X_i]$ for the index already
chosen. Over $L$, its image is divisible by the
minimal polynomial of $x_i$, hence by the original
$P=c_iM_i$ as $c_i$ is a unit of $L$. Primitive
divisibility in Gauss' lemma gives $F\mid H$ already
in $B[X_i]$. The kernel is therefore exactly $(F)$,
and the displayed map is an isomorphism with the
original evaluation map as its defining morphism.
\end{proof}